
\begin{figure}
\centering
\includegraphics[width=1.05\linewidth]{figures/f_equivalence_by_strategy.pdf}
\caption{Artifact equivalence by ecosystem and source-recovery strategy, over artifacts the pipeline could compare (completion rates in \Cref{tab:dataset-and-completion}).}
\label{fig:rq1-equivalence}
\end{figure}

\section{Results}
\label{sec:results}

This section presents the results of our empirical study.

\subsection{RQ1: Ecosystem Verifiability Results}
\label{subsec:results-rq1}

\Cref{fig:rq1-equivalence} reports artifact equivalence across ecosystems and source-recovery strategies, while \Cref{tab:dataset-and-completion} reports pipeline completion and failure causes. We separate these outcomes: \emph{completion} measures whether \toolname recovers the source and rebuilds the artifact, while \emph{equivalence} measures how closely completed rebuilds match the published artifact.

\myparagraph{Internal reproducibility}
We first confirm that \toolname's rebuilds are deterministic by building each package twice.
These rebuilds were byte-identical for 93.0--100.0\% of artifacts across ecosystems and content-identical for 94.8--100.0\%, exceeding reported reproducibility rates for these ecosystems~\cite{benedetti_empirical_2025}. 
% The mismatches reported below therefore reflect genuine source-to-artifact divergence, not noise in our own pipeline.

\myparagraph{Heuristic recovery}
For artifacts without provenance, \toolname recovers the source commit heuristically~(\Cref{subsubsec:stage-1-source-recovery}). Pipeline completion is 74.7\%, 50.9\%, 49.1\%, and 70.9\% for Crates.io, npm, PyPI, and RubyGems, respectively~(\Cref{tab:dataset-and-completion}). The dominant failure is source-commit recovery, a limitation shared by prior rebuild systems and empirical studies that rely on similar heuristics~\cite{imtiaz_dependencies_reviewed_2023,google_oss_rebuild,hassanshahiUnlockingReproducibilityAutomating2025,benedetti_empirical_2025,gao_pyradar_2024,vu_lastpymile_2021,goswami_npm_2020}. Among artifacts that complete, 87.4\% of Crates.io artifacts verify at L1--L3, compared with 52.9\% for npm, 64.6\% for PyPI, and 79.4\% for RubyGems. Crates.io performs best on both completion and equivalence because crate artifacts embed source metadata that directly supports source recovery.

\myparagraph{Provenance-based recovery}
For provenance-bearing packages, \toolname recovers the source commit directly from npm and PyPI attestations, Crates.io Trusted Publishing metadata, or RubyGems attestations. This improves both completion and equivalence. For example, npm completion rises from 50.9\% to 72.4\%, and equivalence among completed artifacts rises from 52.9\% to 73.2\%. RubyGems similarly improves from 70.9\% to 84.5\% completion and from 79.4\% to 85.1\% equivalence, with comparable gains for Crates.io and PyPI~(\Cref{fig:rq1-equivalence}). Once the source commit is known, remaining failures are dominated by build-execution errors and missing build dependencies rather than source recovery.

\myparagraph{Isolating the effect of provenance metadata}
Provenance-bearing packages may also be more popular, active, or mature, so their higher verifiability may not be caused by provenance metadata alone. To isolate the metadata effect, we hold the package set fixed and compare provenance-based recovery with heuristic recovery on the same provenance-bearing packages. Measured as the overall verified fraction of the sample, the direct gain from recorded provenance is small: 0.8, 1.5, 3.6, and $-0.9$ percentage points for Crates.io, npm, PyPI, and RubyGems, respectively. In contrast, the gap between heuristic recovery on provenance-bearing packages and heuristic recovery on the general sample is much larger: 18.2, 24.5, 39.9, and 16.5 points. Thus, most of the verifiability advantage of provenance-bearing packages reflects the build and release practices of adopting projects, not the provenance metadata itself.

\begin{rqanswer}{RQ1 Takeaway}
Verifiability is low and uneven across ecosystems: among completed rebuilds, 52.9--87.4\% verify at L1--L3, and source-commit recovery is the dominant pipeline failure. Provenance improves source recovery and raises apparent verifiability, but holding the package set fixed shows that the recorded metadata contributes at most 3.6 percentage points. Provenance makes source recovery reliable; it does not, by itself, make artifacts verifiable.
\end{rqanswer}

% \Cref{fig:rq1-equivalence} reports pipeline completion and equivalence rates across ecosystems and sample sets, and \cref{fig:rq2-root-causes} summarizes the main causes of pipeline failure.
% Overall, provenance-attested samples achieve higher completion and equivalence rates, while the within-sample comparison shows how much of this improvement comes from using attestation metadata rather than package-population differences.

% For the general ecosystem sample,
% \comment[Oreofe]{"general ecosystem sample", we aren't taking a general ecosystem sample, we could? right now we are doing one full attestation and a sample from the ecosystem, excluding attestation.}
% where \toolname uses heuristic source recovery, \toolname produces equivalence results for \mock{44.0\%}, \mock{49.5\%}, and \mock{67.1\%} of npm, PyPI, and Crates.io artifacts, respectively.
% The main cause of pipeline failure is inability to recover the source commit using the heuristic strategies in \cref{subsubsec:stage-1-source-recovery}.
% These heuristic strategies reflect common source-commit recovery methods used in prior rebuild systems~\cite{imtiaz_dependencies_reviewed_2023, google_oss_rebuild, UnlockReproForOSSHassanshahi} and empirical works~\cite{benedetti_empirical_2025, PyRadar, vu_lastpymile_2021,  goswami_npm_2020},  and this result reflects the limitations of these approaches.
% Over the full sample, equivalence remains low: only \mock{28.4\%}, \mock{32.9\%}, and \mock{60.9\%} of npm, PyPI, and Crates.io artifacts are verified at L1--L3.
% \comment[Oreofe]{I changed the denominator here; it used to say \textit{Among artifacts with equivalence results} (i.e.,\ over compared artifacts), but those rates (64.5 / 66.5 / 90.7\%) don't really read as ``low''. So I switched it to over the full sample (28.4 / 32.9 / 60.9\%) to match the remains low wording. Which denominator do we want? And note PyPI here is just the wheel rate.}
% \PA{I think since our discussion separates pipeline failures from the verifiability results, the verifiability results should be over those the pipeline completed. So we may need to rethink figure 4 to just be over those the pipeline produced results for. Then we introduce a table that captures the pipeline completion rates and reasons why the pipeline failed.}
% Crates.io achieves higher completion and equivalence rates than npm and PyPI, likely because crate artifacts record source metadata that directly supports verification.

% \comment[Oreofe]{We haven't yet gotten a name for the general use of Trusted Pub./Attestation? Provenance-attested applies to npm and pypi but not quite crates.io (and if we do include Ruby too) }
% \PA{Yeah, I'm just realizing that. So we may need to change our terminology. How do you get metadata from Crates trusted publishing and how is it different from vcs info?}
% For the provenance-attested sample using attestation-based recovery, \toolname achieves higher completion and equivalence rates.
% For example, in npm, completion increases by \mock{13.9} percentage points and equivalence increases by \mock{18.8} percentage points.
% These improvements are expected because attestations provide source metadata needed for verification.
% The remaining failures in this setting are primarily due to \mock{build-execution failures} and \mock{missing build dependencies}.

% However, the attested sample may also contain more popular, active, or mature packages, which can independently improve reproducibility.
% \comment[Oreofe]{Does it improve reproducibility or verifiability? Presence of attestation helps us get an almost accurate source commit, workflow, etc., vs. non-attested/TP.}
% To separate these effects, we compare attestation-based and heuristic-based recovery on the same provenance-attested sample.
% Using heuristics instead of attestations leads to \mock{57.6--82.6\%} completion rate and \mock{46.0--79.0\%} equivalence rates across the three ecosystems, which is higher than the general ecosystem sample but still lower than attestation-based results.
% This suggests that only \mock{8--18\%} of the improvement in the attested sample comes from recorded provenance metadata, while \mock{82--94\%} is associated with other package characteristics or development practices in the attestation sample that influence verifiability.

\subsection{RQ2: Root Causes of Non-L1 Verifiability Failures}
\label{subsec:results-rq2}

\newcolumntype{Y}{>{\raggedright\arraybackslash}X}

\begin{table*}[t]
\centering
\scriptsize
\setlength{\tabcolsep}{3.2pt}
\renewcommand{\arraystretch}{1.12}
\caption{Root causes of artifact verifiability failures and the symptoms they produce. We analyzed 134 artifacts, resulting in 185 root cause identifications.}
\label{tab:failure-root-causes}
\begin{tabularx}{\textwidth}{@{}p{0.06\textwidth} p{0.24\textwidth} Y p{0.18\textwidth} c@{}}
\toprule
\textbf{ID} & \textbf{Root Cause} & \textbf{Failure Symptoms} & \textbf{Reconstruction Gap} & \textbf{Count} \\
\midrule
RC01 & Build tool or dependency version differences & Dependency or file version differences in metadata files; version fields of build tools and/or dependencies in manifest or lock files differ; generated files differ because build tool or dependency versions differ. & Environment reconstruction & 47 \\
\midrule
RC02 & CI/CD source code modification & Original artifact contains files or metadata modified by CI/CD workflow steps; package version fields, version-dependent paths, or pre-processed source files differ from the source commit. & Source-state reconstruction & 38 \\
\midrule
RC03 & Generated files or artifacts produced by CI/CD workflow & Generated files or intermediate build outputs such as compiled files or downloaded assets exist in the original artifact but are missing in the rebuilt artifact; files in the original artifact do not exist in the source repository; published artifact represents generated package content. & Build-process reconstruction & 28 \\
\midrule
RC04 & Original timestamp not fixed & Timestamp fields in archives differ. & Build-process reconstruction & 18 \\
\midrule
RC05 & Packaging directory differences & Files in original and rebuilt artifacts were packaged from different temporary build directories. & Build-process reconstruction & 14 \\
\midrule
RC06 & Source revision mismatch & There is a drift between the source commit that produced the artifact and the commit used for rebuilding. & Source-state reconstruction & 11 \\
\midrule
RC07 & Environment-dependent metadata in generated files or file names & Generated files contain local file paths from the build runner environment, or generated file names depend on host paths or directory names. & Environment reconstruction & 6 \\
\midrule
RC08 & Native artifact differences & Native binaries, wheel tags, repair metadata, records, or bundled native files differ. & Environment reconstruction & 5 \\
\midrule
RC09 & Package subpath differences & Files in the rebuilt artifact correspond to a different subdirectory or package in the source repository. & Source-state reconstruction & 4 \\
\midrule
RC10 & Git submodule materialization differences & Files from a Git submodule exist in one artifact but are missing in the other. & Source-state reconstruction & 4 \\
\midrule
RC11 & External VCS source inputs in artifacts & Additional files in original artifacts are VCS-tracked external source payloads, such as Git LFS project files, that were not materialized in the rebuild. & Source-state reconstruction & 4 \\
\midrule
RC12 & Build output contains unstable or unordered inputs & Generated tables, manifests, caches, source files, file names, directories, source maps, or embedded IDs differ across rebuilds. & Build-process reconstruction & 4 \\
\midrule
RC13 & Archive container encoding not fixed & File contents, paths, modes, and mtimes match, but archive SHA differs. & Build-process reconstruction & 1 \\
\midrule
RC14 & VCS line-ending normalization differences & Shared text files differ only by CRLF/LF normalization. & Source-state reconstruction & 1 \\
% \midrule
% RC16 & Build tool configuration differences & Rebuilt artifact contains generated or cache files that are not present in the original artifact. & Build-process reconstruction & 0 \\
% \midrule
% RC17 & Other differences & Rebuilt artifact contains stale directories or files from previous build attempts that do not exist in the original artifact. & Build-process reconstruction & 0 \\
\bottomrule
\end{tabularx}
\end{table*}


% \newcolumntype{Y}{>{\raggedright\arraybackslash}X}

% \begin{table*}
% \centering
% \scriptsize
% \setlength{\tabcolsep}{3.2pt}
% \renewcommand{\arraystretch}{1.12}
% \caption{Root causes of artifact verifiability failures and the symptoms they produce. \mock{Paschal fix the rows in red}}
% \label{tab:failure-root-causes}
% \begin{tabularx}{\textwidth}{@{}p{0.04\textwidth} p{0.24\textwidth} Y p{0.18\textwidth} c@{}}
% \toprule
% \textbf{ID} & \textbf{Root Cause} & \textbf{Failure Symptoms} & \textbf{Reconstruction Gap} & \textbf{Count} \\
% \midrule
% 1 & Build tool or dependency version differences & Version fields of build tools and/or dependencies in manifest or lock files differ. \newline Generated files differ because build tool or dependency versions differ. & Environment reconstruction & 51 \\
% \midrule
% 2 & CI/CD workflow generates intermediate files/artifacts & Generated files or intermediate build outputs such as compiled files or downloaded assets exist in the original artifact but are missing in the rebuilt artifact. & Build-process reconstruction & 28 \\
% \midrule
% 3 & CI/CD workflow modifies source code before packaging & Files in the original artifact contain modifications not present in the source commit. & Source-state reconstruction & 25 \\
% \midrule
% 4 & CI/CD workflow modifies package version & Differences in version fields in manifest files or version-dependent file paths between the original artifact and its source commit. & Source-state reconstruction & 19 \\
% \midrule
% 5 & Original timestamp not fixed & Timestamp fields in the original artifact differ. & Build-process reconstruction & 18 \\
% \midrule
% \mock{6} & Wrong package, directory, or path selected for packaging & Files correspond to a different package, subdirectory, temporary build directory, or path-dependent package selection than the original artifact. & Build-process reconstruction & 17 \\
% \midrule
% \mock{7} & Source revision mismatch & VCS metadata or artifact contents match only after checking out another commit, tag, or temporary release commit. & Source-state reconstruction & 11 \\
% \midrule
% 8 & Environment-dependent paths in generated files or file names & Generated files contain local file paths from the build runner environment, or generated file names depend on host paths or directory names. & Environment reconstruction & 6 \\
% \midrule
% 9 & Native build or wheel repair environment differences & Native binaries, wheel tags, repair metadata, records, or bundled native files differ. & Environment reconstruction & 5 \\
% \midrule
% 10 & Git submodule differences & Files from a Git submodule exist in one artifact but are missing in the other. & Source-state reconstruction & 4 \\
% \midrule
% \mock{11} & Incomplete VCS source inputs & Matching paths contain real payloads in one artifact and pointers or placeholders in the other; tracked files are missing from the rebuilt artifact. & Source-state reconstruction & 4 \\
% \midrule
% \mock{12} & Unstable generated output & Generated tables, manifests, caches, source files, file names, directories, source maps, or embedded IDs differ across rebuilds. & Build-process reconstruction & 4 \\
% \midrule
% \mock{13} & CI/CD source code modification by third-party action & Source files were modified by a CI/CD third-party action before packaging; the original artifact contains modified files not present in the source commit. & Source-state reconstruction & 3 \\
% \midrule
% 14 & Archive container encoding not fixed & File contents, paths, modes, and mtimes match, but archive SHA differs. & Build-process reconstruction & 1 \\
% \midrule
% 15 & VCS line-ending normalization differences & Shared text files differ only by CRLF/LF normalization. & Source-state reconstruction & 1 \\
% \bottomrule
% \end{tabularx}
% \end{table*}


\Cref{tab:failure-root-causes} summarizes the observed root causes, their symptoms, and the reconstruction gaps they expose. For space, we discuss three representative causes.

\myparagraph{Build tool or dependency version differences}
The most common source of divergence was version drift in build tools or build dependencies. 
These differences produced either L3 metadata differences, such as changed lock files, or content differences in files generated by the build toolchain. 
For example, \texttt{tidesurf}~0.2.1 was originally built with Cython~v3.2.3, while our rebuild resolved to Cython~v3.2.6, producing different generated C files.
In such cases, the relevant versions were not pinned in the source manifests, allowing dependency resolution to drift between the original build and the rebuild.

\myparagraph{Generated files produced by CI/CD workflow}
The second most common source of divergence was intermediate files generated by CI/CD workflow steps but not by the canonical build process used in our verifier model. For example, the npm \texttt{lingui-swc}~0.6.0 workflow ran \texttt{yarn build:ts} and \texttt{napi prepublish}, generating files under \texttt{dist/} that were included in the published artifact but absent from our rebuild.
These cases show that artifact verification can fail when build-relevant workflow steps are not captured by registry-exposed metadata or ecosystem-standard build interfaces; prior work similarly finds published packages containing files absent from their source repositories~\cite{vu_lastpymile_2021}.

One might expect that \emph{replaying} the original workflow would close this gap. We evaluated this directly in a pilot: for 105 Crates.io and 105 npm artifacts sampled per outcome bucket, we re-ran each package's real release workflow using \texttt{act}~\cite{nektos_act} and compared the result to the published artifact.\footnote{We excluded other ecosystems after poor results on Crates.io and npm.} Replay improved at most 4\% of artifacts (Crates.io 1/105, npm 4/105) and closed essentially no existing divergences: every npm gain was coverage on a previously failed build rather than a match, and the sole genuine recovery (\texttt{ruckup}~0.9.1, which turned a mismatch into an exact L1 match) succeeded only because the workflow pinned a build-toolchain version our heuristic had guessed wrong. Replay is also fragile: even trivially reproducing artifacts fail to build under \texttt{act} in roughly 9 of 10 cases. Workflow replay therefore does not scale as a verification strategy; the one real gain points instead at build-environment pinning, which we return to in \Cref{subsec:discussion-verifiability-metadata}.

\myparagraph{Source revision mismatch}
Source revision mismatches occur when provenance metadata records a different commit from the one that produced the published artifact.  For example, the Crates.io \texttt{aviutl2}~0.27.0 crate pointed to commit \texttt{7dcd140} and tag \texttt{0.27.0}, while trusted publishing metadata recorded commit \texttt{b2d6bc16}; rebuilding from the recorded commit produced an L1-equivalent artifact. These mismatches often arose from semantic-release-style workflows that modify source metadata, commit the changes, and then publish the artifact; provenance records the triggering commit, not necessarily the final source commit used to produce the artifact.

\begin{rqanswer}{RQ2 Takeaway}
61\% of analyzed divergences were caused by build tool or dependency version drift, workflow-generated intermediate files, or source modifications introduced by CI/CD steps. Improving artifact verifiability requires infrastructure that records build tool versions, source-state changes, and build-relevant workflow steps so that verifiers can accurately reconstruct the source state, environment, and build process used to produce the artifact.
\end{rqanswer}

% \subsection{RQ2: Root Causes of Verifiability Failures}

% \comment[Oreofe]{We can introduce our internal non-determinism here, i.e., we evaluate root causes on cases where we had fully reproducible builds within our pipeline to rule out pipeline-introduced non-determinism, also our reproducibility results are expected because of prior work~\cite{benedetti_empirical_2025}  (and ecosystem efforts?)}
% \PA{Yeah, that makes sense. Really, if our pipeline introduced non-determinism, we should analyze those cases separately, but don't need to analyze them as verifiability failures. Also, should we report reproducibility rate over the whole ecosystem, or we divide cases of non-equivalence into reproducible and not-verifiable and not-reproducible and not-verifiable}

% \begin{figure}[t]
% \centering
% \includegraphics[width=\columnwidth]{figures/rq2_divergence_file_counts.pdf}
% \caption{Number of differing files per mismatched artifact (binned). crates and npm divergences are small; PyPI artifacts differ in many more files.}
% \label{fig:rq2-divergence-file-counts}
% \end{figure}

% \begin{figure}[t]
% \centering
% \includegraphics[width=\columnwidth]{figures/rq2_divergence_file_kinds.pdf}
% \caption{Type of differing files by ecosystem, grouped by execution impact. Source code dominates npm; packaging metadata and other non-executable files dominate PyPI and crates.}
% \label{fig:rq2-divergence-file-kinds}
% \end{figure}

% \Cref{fig:rq2-divergence-file-counts} shows the number of diverging files per artifact, and \cref{fig:rq2-divergence-file-kinds} shows the types of diverging files.
% Most divergences are small: \mock{44\%} of artifacts differ in only one or two files, and only \mock{28\%} differ in more than \mock{10} files.
% Most diverging files are also unlikely to affect execution: \mock{64\%} are metadata, documentation, or other non-source files, while only \mock{36\%} are source files.
% These results support tiered artifact equivalence, which distinguishes non-executable differences from differences that may affect execution.

% \begin{figure}[t]
% \centering
% \includegraphics[width=\columnwidth]{figures/rq2_root_causes.pdf}
% \caption{Root causes of stable equivalence failures by corpus. Packaging-metadata regeneration and Cargo \texttt{vcs\_info} dominate crates; \texttt{npm publish} injection and bundler nondeterminism dominate npm. }
% \label{fig:rq2-root-causes}
% \end{figure}

% \Cref{fig:rq2-root-causes} summarizes the root causes identified in our manual analysis.
% The most common cause is source divergence introduced during the build, where workflow steps or GitHub Actions modify source files before packaging and \toolname does not reproduce those steps.
% Another \mock{XX\%} of divergences arise from build-instruction mismatches: \toolname uses default ecosystem build commands, but some packages rely on custom commands or release-specific workflows.
% A further \mock{XX\%} result from environment-dependent paths, where generated files include absolute build paths in file names or contents. 
% Together, these causes identify gaps in release workflows, build tools, and package metadata that limit artifact verifiability in source-oriented package ecosystems.

% fig:rq3-source-recovery-comparison should be bar chart, with 3 bars per ecosystem for source repository, source commit and package subpath. Each bar should show the proportion of packages we recovered the metadata using heuristics and the proportion that matched what attestation provided.

% fig:rq3-verifiability-comparison is also a bar chart with one bar per ecosystem. The bar should be like a venn chart and show fraction of packages verified by strategy 1, both strategies, strategy 2, no strategy. A question though is whether this 

% \subsection{RQ3: Heuristic vs Attestation-Based Source Recovery}
% % \comment[Oreofe]{Discussions where we should merge this to RQ1 or have as it own}
% % \begin{figure}[t]
% % \centering
% % \includegraphics[width=\columnwidth]{figures/rq3_source_recovery.pdf}
% % \caption{Heuristic source-commit recovery vs.\ the commit recorded by provenance/attestation. ``Recovered but wrong'' = the heuristic returned a commit that disagrees with the attested commit.}
% % \label{fig:rq3-source-recovery-comparison}
% % \end{figure}
% ~\PA{Should we use correct vs incorrect or matched vs not matched, since the heuristics may be more correct. Also, why is n not 4500? It's also surprising how the match rate is not that high (~70\%) but the heuristics performed almost as good as attestation in RQ1.}

% \Cref{fig:rq1-equivalence} compares source metadata recovered by heuristics with metadata from attestations or recorded artifact metadata.
% Across npm and PyPI, heuristic recovery remains incomplete, with the correct source commit recovered for below \mock{75\%} of npm packages.
% Even when heuristics recover metadata, source-commit matches are often low, showing that repository URLs are easier to recover than exact source revisions.
% Crates.io achieves higher recovery and match rates because crate artifacts record source metadata directly.
% These results suggest that ecosystems such as npm and PyPI would benefit from build specifications that record source commits at build time and preserve them in published artifacts.

% \Cref{fig:rq1-equivalence} compares the verifiability rates produced by each recovery method.
% Across ecosystems, attestation-based recovery improves verifiability over heuristic recovery by \mock{18.8}, \mock{42.8}, and \mock{22.0} percentage points for npm, PyPI, and Crates.io.
% \PA{Thinking about this: is the core contribution of attestation the fact it increases the presence of metadata or the quality of metadata?}
% However, some artifacts are verified with heuristic metadata but not with attestation metadata.
% \comment[Oreofe]{Would a reader be interested in a percentage or a fixed number of cases where heuristics were better. An easy way we computed was to see if the commits match; if they don't, we rebuilt the other commit (i.e., the heuristic one}
% \PA{I am not sure I understand. Also, we need to do the analysis in the following paragraph in a more systematic way. We should discuss this more.}
% Our analysis identifies workflow-driven release processes as the most common cause.
% In these cases, GitHub Actions workflows modify source files, create release commits, or tag generated commits during publication.
% The attestation may point to the workflow trigger commit, while tag matching recovers a later release commit that more closely matches the published artifact.
% This finding also exposes a gap in how source commits in provenance attestations are understood by open-source communities: they are often treated as an exact source commit that represents an artifact, but workflow steps can make both the source state and the commit history differ.

% % Six factors is a lot to include in a paper. Based on our result, we can select 2-3 to discuss in the paper, and refer the reader to the remaining in the supplemental materials or discard if they are not interesting.

\begin{figure}[t]
\centering
\includegraphics[width=0.9\linewidth]{figures/rq3_continuous_factors.pdf}
\caption{Verification rate (\% reproduced at L1/L2) by within-ecosystem quartile (Q1 lowest, Q4 highest) of artifact size, downloads, and age, for attested vs.\ non-attested packages across four ecosystems.}
\label{fig:rq3-continuous-factors}
\end{figure}

\subsection{RQ3: Artifact Characteristics and Development Choices}
\label{subsec:results-rq3}

We consider two classes of factors that may explain verifiability outcomes: artifact characteristics, such as size, age, and popularity, and development choices, such as build backend, native code, and repository layout. \Cref{fig:rq3-continuous-factors} shows that artifact characteristics have limited explanatory power: verifiability declines only slightly with artifact size, shows little relationship with popularity, and the apparent effect of release age largely disappears once attestation is controlled for. In contrast, development choices have a much stronger effect. Across ecosystems, verifiability is primarily shaped by how deterministically the build and packaging process reconstructs the published artifact.

\myparagraph{Packaging determinism matters more than closeness to source}
Crates.io performs consistently well because crates are published as source tarballs through a standardized packaging process, with remaining mismatches mostly caused by embedded build timestamps~(\Cref{subsec:results-rq2}). However, being close to source is not sufficient. Package artifacts can diverge from their repositories even without compilation, because packaging can rewrite metadata, normalize files, or apply preprocessing~\cite{vu_lastpymile_2021}. npm illustrates this effect: \texttt{pack-only} packages verify less often than bundled or transpiled packages, even though the latter transform the source more heavily. The difference is packaging determinism: stable generated outputs can be easier to reproduce than source-like artifacts produced by nondeterministic packaging steps.

\myparagraph{PyPI has high-verifiability backends \& low-verifiability paths}
PyPI shows the widest spread across build backends. At the high end, \texttt{hatchling} reproduces many wheels at L1--L2, requiring little normalization. Other backends, including \texttt{pdm-backend}, \texttt{poetry-core}, and \texttt{setuptools}, reach comparable total verification rates mainly through L3 normalization, indicating that their outputs often differ in metadata, timestamps, or ordering rather than executable content. At the low end, wheels built through legacy \texttt{setup.py bdist\_wheel} paths or packages without a declared backend rarely verify. This spread helps explain why PyPI contains some highly verifiable build configurations while remaining difficult to verify overall.

\myparagraph{Native compilation reduces verifiability}
Native compilation substantially lowers verification rates in ecosystems that ship built artifacts. On PyPI, wheels with compiled extensions, such as those built with \texttt{maturin} or \texttt{scikit-build-core}, verify far less often than pure-Python wheels; on npm, \texttt{node-gyp} modules similarly sit near the bottom of the ecosystem. These packages depend on platform-specific compilers, flags, and toolchain behavior, introducing variation that package-level normalization cannot remove~\cite{malkaReproducibilityBuildEnvironments2024,randrianainaOptionsMatterDocumenting2024}. Native packaging therefore remains a major barrier to independent artifact verification.

\myparagraph{Monorepos mainly affect source recovery}
Repository layout has little effect once the correct package directory is recovered. PyPI packages built from monorepo subdirectories verify at the same rate as root-level packages, suggesting that monorepos primarily create a source-recovery challenge rather than a rebuild or comparison challenge. Their impact therefore appears in pipeline completion failures~(\Cref{tab:dataset-and-completion}), not in equivalence outcomes among completed rebuilds.

\begin{figure}[t]
\centering
\includegraphics[width=\linewidth]{figures/rq3_dev_categorical.pdf}
\caption{Packages verified (\%) by build backend or packaging strategy, grouped by ecosystem and split into L1--L2 and L3 matches.}
\label{fig:rq3-dev-categorical}
\end{figure}

\begin{rqanswer}{RQ3 Takeaway}
Verifiability is shaped more by how an artifact is built than by what the artifact is. Size, age, and popularity have limited effects, while deterministic packaging, declarative build backends, and avoidance of native compilation strongly improve verification outcomes. Monorepos mainly affect source recovery rather than rebuild equivalence, so the most important levers for improving verifiability are build- and packaging-side changes.
\end{rqanswer}

\subsection{RQ4: \toolname vs.\ OSS-Rebuild}
\label{subsec:results-rq4}

Across the 600 artifacts, \toolname verifies 331 (55\%) versus 158 (26\%) for
OSS-Rebuild---more than twice as many. \Cref{tab:rq4-comparison} provides the details. 
\toolname attains the higher verification rate in four of the six strata, ties in one (PyPI unsupported), and is exceeded only on supported npm packages.

\myparagraph{\toolname generalizes to the long tail; OSS-Rebuild does not}
The starkest gap is the unsupported long tail: on unsupported crates.io, OSS-Rebuild verifies \emph{none} (0\%) against \toolname's 74\%, and 74\%, 22\%, and 6\% of unsupported crates, npm, and PyPI artifacts respectively are verified by \toolname alone. OSS-Rebuild depends on curated, per-package build definitions, so outside its supported set it often has no recipe to run---most visibly on crates.io, where it rebuilt no unsupported package. \toolname instead recovers the published source commit and rebuilds from it, transferring to arbitrary packages without curation.
\begin{table}[t]
  \centering
  \caption{Verification by \toolname{} and OSS-Rebuild on 600 artifacts
  (100 per ecosystem/stratum).}
  \label{tab:rq4-comparison}
  \scriptsize
  % \setlength{\tabcolsep}{2.2pt}
  % \renewcommand{\arraystretch}{0.92}
  \resizebox{0.9\columnwidth}{!}{%
  \begin{tabular}{@{}llrrrr@{}}
    \toprule
    Eco. & Strat. & \toolname-only & OSS-only & Both & None \\
    \midrule
    crates & supp.    & 61\% & 2\%  & 25\% & 12\% \\
    crates & un-supp. & 74\% & 0\%  & 0\%  & 26\% \\
    npm    & supp.    & 6\%  & 13\% & 33\% & 48\% \\
    npm    & un-supp. & 22\% & 13\% & 21\% & 44\% \\
    PyPI   & supp.    & 44\% & 4\%  & 22\% & 30\% \\
    PyPI   & un-supp. & 4\%  & 6\%  & 19\% & 71\% \\
    \midrule
    \multicolumn{2}{@{}l}{Overall} & 35\% & 6\% & 20\% & 39\% \\
    \bottomrule
  \end{tabular}%
  }
  \vspace{-0.5em}
  % \begin{flushleft}
  % \scriptsize
  % % \toolname{}/OSS: total share verified by each platform (L1--L3); Both: verified
  % % by both; \toolname-only/OSS-only: verified by exactly one; None: verified by
  % % neither.
  % \end{flushleft}
\end{table}

% \OS{AVP-only, OSS-only, Both and None}
\myparagraph{\toolname is competitive even where OSS-Rebuild is tuned}
Even on OSS-Rebuild's supported strata, \toolname verifies more crates.io (86\% vs.\ 27\%) and PyPI (66\% vs.\ 26\%) artifacts; OSS-Rebuild leads only on supported npm (46\% vs.\ 39\%), where its hand-maintained definitions reproduce bundled and transpiled JavaScript that \toolname's generic \texttt{npm pack} does not. The platforms are complementary, not redundant (\Cref{tab:rq4-comparison}): on supported crates.io they agree on only 25\% of artifacts while \toolname uniquely verifies a further 61\%, and OSS-Rebuild uniquely verifies only a small share anywhere (0--13\%). A non-trivial 12--71\% is verified by neither---native compilation, undeclared build backends, and unrecoverable source.

\begin{rqanswer}{RQ4 Takeaway}
\toolname verifies about twice as many artifacts as OSS-Rebuild overall and is the only platform that scales beyond a curated package set. It trades some per-package fidelity on heavily transformed artifacts (e.g., supported npm) for broad coverage that needs no per-package build definitions, making generality the decisive advantage for assessing artifact authenticity at scale.
\end{rqanswer}

% \comment[Oreofe]{Note: our artifact-characteristic fig and development-choices are referenced here but we haven't built them yet. On the data: we have size and age for all three ecosystems, but downloads only for PyPI, and monorepo/native aren't columns we 
% have; we'd have to derive them (native we can partly get from \texttt{build\_backend}=maturin/scikit-build-core, which is already in \texttt{fig:rq4-build-system-impact}). So do we build characteristics as just size+age quartiles, and is dev-choices worth 
% the extra derivation work? Only the build-system panel exists right now.}
% \PA{Is there a uniform popularity measurement we can use that is present across all ecosystems? Maybe number of github stars? Is it possible to get the monorepo status and native status in any way? We can approximate monorepo by checking whether the package directory recovered was the root directory or a subdirectory. FOr native, I think a more sound approach will be to unpack the original artifact and check for native file extensions. Alternative, we can discuss if these are worth the effort or if there are other development practices we can measure or if we just leave it at build system choices}
% \Cref{fig:rq4-artifact-characteristics} shows the relationship between artifact characteristics and verifiability.
% Smaller artifacts have higher verifiability rates, consistent with the expectation that larger artifacts provide more opportunities for content divergence.
% Popularity shows only minor differences in verifiability.
% Surprisingly, newer artifacts have lower verifiability rates, suggesting that recent package releases may use more complex or less deterministic build processes.
% These relationships vary by ecosystem and are more pronounced in PyPI than in Crates.io, likely because Crates.io has higher overall verifiability rates.

% \Cref{fig:rq4-development-choices} shows the relationship between development choices and verifiability.
% Monorepos and native extensions are associated with lower verifiability rates, likely because they require more complex source recovery, build configuration, or platform-specific dependencies.
% \Cref{fig:rq4-build-system-impact} further shows that default or widely used build tools tend to have higher verifiability rates than less common alternatives.

% \begin{figure}[t]
% \centering
% \includegraphics[width=\columnwidth]{figures/rq4_build_system.pdf}
% \caption{PyPI wheel verifiability (L1/L2) by build backend. Modern declared backends (hatchling, pdm, setuptools) verify well; native-extension backends (maturin, scikit-build-core) and the legacy \texttt{bdist\_wheel} path do not.}
% \label{fig:rq4-build-system-impact}
% \end{figure}
% These results suggest that artifact verification tools must account for ecosystem-specific build systems, monorepo layouts, and native-code packaging to improve verifiability in package ecosystems.

% \begin{table}[ht!]
% \footnotesize % Steps down one size from \small to fit IEEE columns better
% \centering
% \caption{Ecosystem-wide package and repository statistics.}
% \label{tab:ecosystem-stats}
% \setlength{\tabcolsep}{4pt} % Reduces white space between columns (default is usually 6pt)
% \renewcommand{\arraystretch}{1.1}
% \begin{tabular}{@{} l r r r @{}} % @{} removes outer edge padding to maximize space
% \toprule
% & \textbf{PyPI} & \textbf{npm} & \textbf{crates.io} \\
% \midrule
% Total unique packages & 748,031 & 3,991,783 & 261,674 \\
% Has attested prov.    & 49,436 (6.6\%) & 104,322 (2.6\%) & 5,206 (2.0\%) \\
% Has any repo URL      & --- & 2,451,966 (61.4\%) & 223,261 (85.3\%) \\
% Has GitHub URL        & 431,168 (57.6\%) & 2,200,440 (55.1\%) & 211,490 (80.8\%) \\
% No valid repo URL     & 316,863 (42.4\%) & 1,539,817 (38.6\%) & 38,413 (14.7\%) \\
% \bottomrule
% \end{tabular}
% \end{table}

% \comment[Oreofe]{Across the 748,031 unique PyPI packages in our dataset, we identified a valid GitHub repository URL for 431,168 packages (57.6\%). The remaining 316,863 packages (42.4\%) did not have a valid GitHub URL that could be identified using our extraction and validation procedure. This result is consistent with prior ecosystem-scale studies. Tsakpinis and Pretschner reported that the proportion of PyPI libraries with valid GitHub repository URLs ranged from 58.0\% to 73.8\%, depending on package importance~\cite{PypiNpmGHAccessibility}. More recent work found that 65.8\% of PyPI libraries linked to a source code repository, that GitHub accounted for 97.2\% of repository links, and that 12.9\% of GitHub links were outdated~\cite{tsakpinis2026investigatingnotablemetadatapractices}. Gao et al. similarly found that metadata-based tools could retrieve repository information for at most 63.1\% of PyPI packages after accounting for URL redirection~\cite{PyRadar}. Collectively, these results indicate that missing, outdated, or otherwise unusable repository metadata remains a substantial obstacle to ecosystem-scale artifact verification.
% }

% \subsection{RQ1: What Proportion of Software Packages Can Be Independently Verified?}

% We first measure the proportion of packages that \toolname can independently verify against their published registry artifacts. For each ecosystem, we report the strongest equivalence level achieved for each package across the applicable metadata-recovery strategies and reconstructed build configurations. Following the definitions in Section~\ref{sec:equivalence-levels}, we treat L1 byte-for-byte matches and L2 extracted-content matches as strict verification outcomes, while reporting weaker equivalence levels separately. 

% We report two complementary measures. First, the \emph{comparison coverage} captures the proportion of sampled packages for which \toolname can recover sufficient metadata, successfully rebuild the package, and obtain an artifact suitable for comparison. Second,
% % might want to think of a better name
% the \emph{full verification yield}
% captures the proportion of all sampled packages that \toolname verifies at L1 or L2. 

% \begin{comment}
%     - 
%     - Histogram (Bar chart) of files that differ in each ecosystem
% \end{comment}

% \comment[Oreofe]{Add Reproducibility Results to table; Portion that Not Verifiable but Build was reproducible}
% \begin{table*}[!t]
%     \centering
%     \caption{\toolname outcomes by ecosystem and artifact format}
%     \label{table:rq1}
%     \scriptsize
%     \begin{tabular}{lllrrrrrrrr}
%         \toprule
%         \textbf{Ecosystem} &
%         \textbf{Corpus} &
%         \textbf{Artifact} &
%         \textbf{Sampled} &
%         \textbf{Compared} &
%         \textbf{L1} &
%         \textbf{L2} &
%         \textbf{L3} &
%         \textbf{L4} &
%         \textbf{L5+} &
%         \textbf{Strict E2E} \\
%         \midrule
%         PyPI &
%         Attested &
%         Wheel &
%         4{,}500 &
%         3{,}885 (86.3\%) &
%         1{,}473 (37.9\%) &
%         1{,}510 (38.9\%) &
%         0 (0.0\%) &
%         456 (11.7\%) &
%         446 (11.5\%) &
%         2{,}983 (66.3\%) \\
        
%         PyPI &
%         Attested &
%         Sdist &
%         4{,}500 &
%         3{,}718 (82.6\%) &
%         1{,}214 (32.7\%) &
%         1{,}827 (49.1\%) &
%         0 (0.0\%) &
%         258 (6.9\%) &
%         419 (11.3\%) &
%         3{,}041 (67.6\%) \\
        
%         PyPI &
%         Non-attested &
%         Wheel &
%         4{,}500 &
%         2{,}204 (49.0\%) &
%         246 (11.2\%) &
%         271 (12.3\%) &
%         0 (0.0\%) &
%         897 (40.7\%) &
%         790 (35.8\%) &
%         517 (11.5\%) \\

%         PyPI &
%         Non-attested &
%         Sdist &
%         4{,}500 &
%         2{,}591 (57.6\%) &
%         124 (4.8\%) &
%         386 (14.9\%) &
%         0 (0.0\%) &
%         316 (12.2\%) &
%         1{,}765 (68.1\%) &
%         510 (11.3\%) \\
%         \midrule
%         npm &
%         Attested &
%         Tarball &
%         4{,}500 &
%         2{,}640 (58.7\%) &
%         1{,}051 (39.8\%) &
%         535 (20.3\%) &
%         0 (0.0\%) &
%         233 (8.8\%) &
%         821 (31.1\%) &
%         1{,}586 (35.2\%) \\

%         npm &
%         Non-attested &
%         Tarball &
%         4{,}500 &
%         1{,}979 (44.0\%) &
%         148 (7.5\%) &
%         520 (26.3\%) &
%         0 (0.0\%) &
%         77 (3.9\%) &
%         1{,}234 (62.4\%) &
%         668 (14.8\%) \\
%         \midrule
%         crates.io &
%         Trusted Publishing &
%         Crate &
%         5{,}206 &
%         4{,}390 (84.3\%) &
%         2{,}306 (52.5\%) &
%         181 (4.1\%) &
%         0 (0.0\%) &
%         1{,}547 (35.2\%) &
%         356 (8.1\%) &
%         2{,}487 (47.8\%) \\

%         crates.io &
%         VCS metadata &
%         Crate &
%         4{,}500 &
%         3{,}020 (67.1\%) &
%         944 (31.3\%) &
%         221 (7.3\%) &
%         0 (0.0\%) &
%         1{,}415 (46.9\%) &
%         440 (14.6\%) &
%         1{,}165 (25.9\%) \\
%         \bottomrule
%     \end{tabular}
% \end{table*}

% \paragraph{PyPI.} \toolname verifies a substantial fraction of packages in the PyPI attested corpus. For wheels, \toolname reaches artifact comparisoon for 3{,}885 of the 4{,}500 sampled packages (86.3\%). Among these packages, 1{,}473 wheels (37.9\%) match the published artifact byte-for-byte, while an additional 1{,}510 wheels (38.9\%) contain identical extracted file sets and contents despite archive-level differences.
% Thus, \toolname strictly verifies 2{,}983 attested packages at L1 or L2, corresponding to 76.8\% of compared wheels and 66.3\% of the full attested corpus.

% The results are similar for source distributions.
% \toolname reaches sdist comparison for 3{,}718 attested packages (82.6\%) and strictly verifies 3{,}041 packages at L1 or L2.
% This corresponds to 81.8\% of compared sdists and 67.6\% of the full attested corpus.
% For both artifact formats, L2 matches are common: 38.9\% of compared wheels and 49.1\% of compared sdists have equivalent payloads despite archive-level differences.
% Consequently, byte-for-byte equality alone would substantially understate the proportion of artifacts whose packaged contents can be independently reconstructed.

% Verification is substantially more difficult in the non-attested PyPI corpus.
% \toolname reaches wheel comparison for 2{,}204 packages (49.0\%) and strictly verifies 517 packages, corresponding to 23.5\% of compared wheels and 11.5\% of the full corpus.
% For sdists, \toolname reaches comparison for 2{,}591 packages (57.6\%) and strictly verifies 510 packages, corresponding to 19.7\% of compared sdists and 11.3\% of the full corpus.
% Residual mismatches also differ by artifact format.
% Among compared non-attested wheels, 40.7\% differ only in generated packaging metadata at L4.
% In contrast, 68.1\% of compared non-attested sdists remain non-equivalent at L5, most commonly because the published and rebuilt archives contain different file sets.
% We investigate the causes of these mismatches in RQ2 (See Section~\ref{sec:rq2}).

% \paragraph{npm.} For npm, \toolname obtains a comparable registry tarball for 2{,}640 packages in the attested corpus.
% Among these packages, 1{,}051 tarballs (39.8\%) match the published artifact byte-for-byte at L1, while an additional 535 packages (20.3\%) match at the extracted-content level at L2.
% Thus, \toolname strictly verifies 1{,}586 packages at L1 or L2, corresponding to 60.1\% of packages that reach artifact comparison and \emph{[X]\%} of the full attested corpus.
% A further 233 packages (8.8\%) differ only in generated packaging metadata at L4.
% The remaining 821 packages (31.1\%) remain non-equivalent at L5, most commonly because the rebuilt and published tarballs contain different file sets.
% For the non-attested npm corpus, \toolname obtains a comparable tarball for 1{,}979 packages.
% Of these, 148 packages (7.5\%) match byte-for-byte and an additional 520 packages (26.3\%) match at the extracted-content level.
% Together, 668 packages satisfy strict L1 or L2 equivalence, corresponding to 33.8\% of compared packages and \emph{[Y]\%} of the full non-attested corpus.
% Only 77 packages (3.9\%) stop at a metadata-only mismatch.
% The remaining 1{,}234 packages (62.4\%) remain non-equivalent at L5: 576 packages (29.1\%) exhibit file-set-only differences, 304 packages (15.4\%) exhibit content-only differences, and 354 packages (17.9\%) exhibit both file-set and content differences.

% \paragraph{crates.io.} For crates.io, we analyze the complete population of 5{,}206 crates for which \toolname recovers source-linkage evidence through Trusted Publishing.
% \toolname successfully rebuilds 4{,}390 crates (84.3\%).
% Among these buildable crates, 2{,}306 (52.5\%) produce byte-for-byte identical archives at L1, while an additional 181 (4.1\%) match at the extracted-content level at L2.
% Thus, \toolname strictly verifies 2{,}487 crates, corresponding to 56.7\% of buildable crates and 47.8\% of the complete Trusted Publishing population.

% A further 1{,}547 buildable crates (35.2\%) differ only in generated Cargo metadata at L4.
% The most frequently differing metadata file is \textit{Cargo.lock}, followed by \textit{.cargo\_vcs\_info.json} and \textit{Cargo.toml}.
% If these metadata-only outcomes are reported separately as a weaker equivalence level, \toolname reaches L1--L4 for 91.9\% of buildable crates and 77.5\% of the complete population.
% The remaining 356 buildable crates (8.1\%) exhibit file-set differences, content differences, or both.
% \paragraph{Cross-ecosystem comparison.}

% We isolate the effect of the metadata-recovery strategy using paired comparisons in RQ3.

% \begin{rqanswer}{RQ1 Takeaway}
% We
% \end{rqanswer}

% \subsection{RQ2: Root Causes of Verifiability Failures}
% \toolname fails to strictly verify a package for one of two reasons. First, \emph{pipeline faliures} prevent
% the package from reaching artifact comparison at all: \toolname cannot recover a source revision, cannot reconstruct
% a working build environment, the rebuild itself fails, or two independent rebuilds disagree. Second, \emph{equivalence faliures}
% occur when a package rebuilds successfully but the rebuilt artifact does not match the published
% artifact at L1 or L2. Throughout, \toolname's double-rebuild check lets us separate nondeterminism introduced by our own pipeline from
% stable divergence attributable to the package, its original build, or its publication path (Section~\ref{sec:rq2-repeatability}).
% Table~\ref{tab:rq2-pipeline} decomposes the pipeline failures for each corpus, and Table~\ref{tab:rq2-rootcauses} presents our root-cause taxonomy for
% equivalence failures. We analyze each failure class in turn and then compare our taxonomy against failure causes reported in prior
% reproducible-build studies.

% \subsubsection{Internal Build Repeatability}
% \label{sec:rq2-repeatability}
% \begin{table}[t]
% \caption{Internal build repeatability: each build configuration is
% executed twice and the two rebuilds are compared against each
% other. Counts are over build pairs, not packages.}
% \label{tab:rq2-repeat}
% \centering
% \footnotesize
% % \setlength{\tabcolsep}{3pt}
% % \begin{tabular}{llrrrr}
% \begin{tabular}{@{} l l r r r  @{}} % @{} removes outer edge padding to maximize space

% \toprule
% Corpus & Artifact & Pairs & Byte-ident. & Content-ident.\\
% \midrule
% PyPI Att.       & Wheel   & 13,295 & 12,525 (94.2\%) & 0 (0.0\%)  \\
% PyPI Att.       & Sdist   & 12,708 & 6,301 (49.6\%)  & 6,129 (48.2\%) \\
% PyPI Non-att.   & Wheel   & 4,253  & 4,167 (98.0\%)  & 0 (0.0\%)  \\
% PyPI Non-att.   & Sdist   & 5,024  & 853 (17.0\%)    & 4,136 (82.3\%) \\
% npm Att.        & Tarball & 2,981  & 2,816 (94.5\%)  & 42 (1.4\%) \\
% npm Non-att.    & Tarball & 1,949  & 1,865 (95.7\%)  & 8 (0.4\%) \\
% crates.io TP    & Crate   & 4,389  & 4,389 (100.0\%) & 0 (0.0\%) \\
% crates.io VCS   & Crate   & 3,001  & 3,001 (100.0\%) & 0 (0.0\%) \\
% crates.io TP-VCS\footnotemark & Crate & 4,173 & 4,173 (100.0\%) & 0 (0.0\%) \\
% \bottomrule
% \end{tabular}
% \end{table}
% \footnotetext{Trusted Publishing crates rebuilt via the
% \texttt{.cargo\_vcs\_info.json} recovery path.}

% Before attributing mismatches to root causes, we must rule out nondeterminism introduced by \toolname itself. \toolname executes every build
% configuration twice in the same reconstructed environment and compares the two rebuilds against each other, mirroring the L1/L2
% distinction: a build pair is \emph{byte-repeatable} if the two rebuilt artifacts are byte-for-byte identical, and \emph{content-repeatable} if their extracted 
% file sets and contents are identical despite archive-level differences. Table~\ref{tab:rq2-repeat} reports repeatability over all executed
% build pairs.

% \paragraph{Cargo is fully deterministic.}
% crates.io exhibits zero mismatching pairs across all 11,563 build pairs: given a pinned Rust toolchain and timestamp strategy,
% \texttt{cargo package} produces byte-identical archives on every run. Consequently, every crates.io mismatch against a published
% artifact in the analysis below reflects stable divergence: wrong recovered metadata, environment differences, or genuine divergence
% between public source and registry artifact, and never rebuild noise. Benedetti et al. confirm the same as they find \dots~\cite{benedetti_empirical_2025}

% \paragraph{Byte-level repeatability caps L1 verification.}
% For sdists, byte-repeatability is low (49.6\% attested, 17.0\% non-attested) while content-repeatability is higher (97.8\% and 99.3\% cumulatively): nearly all internal byte
% divergence is tar and gzip archive metadata that varies across consecutive runs. If \toolname's own back-to-back rebuilds differ at the
% archive level, the rebuild cannot match the published artifact byte-for-byte either. The attested-sdist L1 rate of 32.7\% (Table~\ref{tab:rq1}) must therefore be read against an internal
% ceiling of 49.6\%, and content-level (L2) equivalence is the appropriate strict tier for sdists. Wheels and crates show no such gap. As a consistency check, every package verified at L1 is also
% byte-repeatable (\mock{zero} exceptions).
% % TODO(analysis): re-run the L1 => byte-repeatable check at scale

% \paragraph{Conditioning the failure analysis.}
% These results license the root-cause analysis that follows. For a content-repeatable package, a mismatch against the published artifact is \emph{stable}: it cannot be attributed to stochastic
% noise in our rebuild pipeline. Of the 11,470 artifact comparisons that fail strict verification (L4 and L5 in Table~\ref{tab:rq1}), \mock{N} (\mock{$\sim$97\%}) come from content-repeatable builds; 
% we restrict the root-cause taxonomy to this population. The remaining \mock{n} packages have no content-repeatable build configuration even after eliminating
% harness-introduced nondeterminism (see Internal nondeterminism, Section~\ref{sec:rq2-pipeline}); these are unverifiable by \emph{any} rebuild-based verifier under our build approach,
% regardless of recovery strategy. We count them under the Nondet.\ column of Table~\ref{tab:rq2-pipeline} and report them separately rather than folding them into the taxonomy.